\section*{1.1 Our Result}
We state our result formally now. Without loss of generality, we will assume that a depth four circuit is a $\Sigma\Pi\Sigma\Pi$ circuit, i.e., the circuit has a $+$-gate on top followed by second layer of $\times$-gates, then a third layer of $+$-gates and finally a bottom layer of $\times$-gates.
\begin{authorclaim}{Theorem1 (Lower bound for depth four circuits)}
Over any field of characteristic zero\footnote{The lower bound holds even if the characteristic is sufficiently large (see Section4).}, there exists a family of mulitilinear polynomials $\{f_n\}_{n\ge1}$ in VNP, where $f_n$ is a polynomial in $\Theta(n)$ variables and of degree $\Theta(n)$ such that any depth four circuit computing $f_n$ has $\Omega(n^{2.5}/(\log n)^6)$ many wires/edges and $\Omega(n^{1.5}/(\log n)^4)$ many gates.
\end{authorclaim}
\paragraph{A word about the polynomial family}
The polynomial family $\{f_n\}_{n\ge1}$ is a variant of the Nisan-Wigderson design polynomial family, which has been used in proving several other previous lower bounds [49,51–53,57,58,73,85]. The $n$-th member of the family, i.e., $f_n$ is a polynomial in $\mathbf x=\{x_1,\ldots,x_{3m}\}$ and $\mathbf y=\{y_1,\ldots,y_{3m}\}$ variables, where $m$ is an integer in $[n/2,2n]$. The degree of the polynomial in $\mathbf y$ variables is $\deg_{\mathbf y}(f_n)=m$, while its degree in $\mathbf x$ variables is $\deg_{\mathbf x}(f_n)=d_x=\Theta(\sqrt m/\ln m)$. Informally, $f_n$ contains multiple 'copies' of the design polynomial in different subsets of the $\mathbf x$ variables, while the $\mathbf y$ variables are used as 'prefixes' to uniquely identify each such copy. While the reason for having multiple copies is similar to [53], as we shall see in the next section, handling them is a little trickier in our case. Note that because of the way we have defined $m$ and $d_x$, proving that any depth four circuit computing $f_n$ has $\Omega(m^2d_x/(\ln m)^5)$ many edges and $\Omega(md_x/(\ln m)^3)$ many gates would establish the theorem.\footnote{We use log base $e$ in the proof rather than log base2 as it simplifies the analysis.} The exact description of $f_n$ is given in Section4.
